% keycheat.tex — the rendering engine. Content files (sections/*.tex)
% and layout.tex should only ever call the commands defined here; they
% should never need to know how a keycap or a group box is actually drawn.
%
% Public API:
%   \key{Label}                     -> one keycap, usable inline anywhere
%   \combo{Mod,Shift,E}             -> a run of keycaps joined by "+"
%   \sr{Mod,Shift,E}{description}   -> one table row: combo + description
%   \srx{<raw left cell>}{description} -> one table row, custom left cell
%   \group{Title}{Width}{Rows}[Note] -> a titled panel; Rows is a tabular body
%     (built from \sr / \srx lines), Note is an optional muted footer line
%
% Styling knobs (colours, fonts) live in palette.tex / config.tex, not here.

\tikzset{
  shadow xshift=0pt,
  shadow yshift=-2bp,
  keycapstyle/.style={
    rounded corners=3pt,
    draw=borderKeycap,
    line width=0.6pt,
    fill=bgKeycap,
    minimum height=1.8em,
    minimum width=1.8em,
    inner sep=4pt,
    font=\keyfont\small,
    text=textPrimary,
    align=center,
  },
  groupbox/.style={
    rounded corners=7pt,
    draw=borderGroup,
    line width=1pt,
    fill=bgGroup,
    inner sep=14bp,
    align=left,
    drop shadow={opacity=0.35,fill=black},
  },
}

\newcommand{\key}[1]{\tikz[baseline=(k.base)]{\node[keycapstyle](k){#1};}}

\newcommand{\keyplus}{\,{\footnotesize\color{textSecondary}+}\,}

\newcommand{\combo}[1]{%
  \foreach \k [count=\i] in {#1}{%
    \ifnum\i>1 \keyplus\fi
    \key{\k}%
  }%
}

\newcommand{\sr}[2]{\combo{#1} & #2 \\}
\newcommand{\srx}[2]{#1 & #2 \\}

\NewDocumentCommand{\group}{mmmO{}}{%
  \begin{minipage}[t]{#2 bp}
    {\keyfont\bfseries\color{accent}\fontsize{15}{15}\selectfont #1}\\[4bp]
    {\color{borderGroup!60!bgGroup}\rule{\linewidth}{0.8bp}}\\[7bp]
    \renewcommand{\arraystretch}{1.45}
    \setlength{\tabcolsep}{0pt}
    \begin{tabular}{@{}r@{\hspace{9bp}}l@{}}
      \small\color{textPrimary}
      #3
    \end{tabular}
    \IfNoValueTF{#4}{}{\\[7bp]{\fontsize{8}{9}\selectfont\itshape\color{textSecondary}#4}}
  \end{minipage}%
}
